%%%%%%%%%%  scan idx 69  |  folio 56  %%%%%%%%%%
\SGAPrintPage{56}%
\sgaStableTarget{sga6.e00.eq.2-30}\begin{equation}\tag{2.30}
 \underline K(Y)\longrightarrow\underline K(X)\longrightarrow\underline K(U)\longrightarrow 0
\end{equation}
\emph{is exact.}

N.B. Compare with proposition 2.3.

The fact that $\underline K(X)\longrightarrow\underline K(U)$ is surjective follows immediately from proposition 2.5 (or also from proposition 2.6). It remains to prove that, if $x\in\underline K(X)$ has zero restriction to $\underline K(U)$, it is in the image of $\underline K(Y)$. For this, we return to the explicit definition of $\underline K(X)$ and $\underline K(U)$: we have two coherent algebraic sheaves $F$ and $G$ on $X$ with $x=\gamma_X(F)-\gamma_X(G)$ and $\gamma_U(F)-\gamma_U(G)=0$; this last relation means that we can find coherent sheaves $P'$ and $Q'$ on $U$, equipped with filtrations whose factors are pairwise isomorphic on $U$, and such that we have an isomorphism on $U$:
\[
 P=F+Q'\longrightarrow Q=G+P'.
\]
By proposition 2.5, the sheaves $P'$ and $Q'$ are restrictions to $U$ of coherent sheaves on $X$ (denoted by the same letters); likewise, their filtrations come from filtrations on the corresponding sheaves on $X$, by lemma 2.11 below. We then have:
\[
 \gamma_X(F)-\gamma_X(G)=(\gamma_X(P)-\gamma_X(Q))+(\gamma_X(P')-\gamma_X(Q')).
\]
Introducing the factors $P'_i$ and $Q'_j$ of the filtrations of $P'$ and $Q'$ on $X$, the last term is a sum of terms $\gamma_X(P'_i)-\gamma_X(Q'_i)$, where $P'_i|U\simeq Q'_i|U$. We are thus reduced to proving the following particular case: if $R$ and $S$ are coherent sheaves on $X$ such that $R|U\simeq S|U$, then $\gamma_X(R)-\gamma_X(S)$ is in the image of $\underline K(Y)$. Now let $T$ be the graph of the given isomorphism from $R|U$ onto $S|U$; it is a coherent subsheaf of $(R\times S)|U$, and therefore, by lemma 2.11, it is the restriction to $U$ of a coherent subsheaf,
%%%%%%%%%%  scan idx 70  |  folio 57  %%%%%%%%%%
\SGAPrintPage{57}%
again
denoted $T$, of $R\times S$. We then have $\gamma_X(R)-\gamma_X(S)=(\gamma_X(R)-\gamma_X(T))-(\gamma_X(S)-\gamma_X(T))$, and it suffices to prove, for example, that the first term of the second member is in the image of $\underline K(Y)$. Now the projection $u:T\longrightarrow R$ is an isomorphism over $U$, so $\gamma_X(R)-\gamma_X(T)=\gamma_X(\operatorname{Coker}u)-\gamma_X(\operatorname{Ker}u)$, and $\operatorname{Ker}u$ and $\operatorname{Coker}u$ have their support in $Y$.

It therefore remains to prove the following lemma:

\sgaStableTarget{sga6.e00.lem.2-11}\textbf{Lemma 2.11.} \emph{Let $X$ be an algebraic space, $F$ a coherent algebraic sheaf on $X$, $U$ an open subset of $X$, and $G$ a coherent subsheaf of $F|U$. Then $G$ is the restriction to $U$ of a coherent subsheaf of $F$.}

Let $\overline G$ be the subsheaf of $F$ whose sections on the open subset $V$ of $X$ are the sections of $F$ on $V$ whose restriction to $U\cap V$ is a section of $G$. Everything reduces to proving that this sheaf $\overline G$ is coherent. For this, we can assume $X$ and $U$ affine (for $U$ is a union of affine open subsets $U_i$, and $\overline G$ is then the intersection of the analogous sheaves $\overline{G|U_i}$). Then the coordinate ring of $U$ is a ring of fractions $A_S$ of the coordinate ring $A$ of $X$, and we verify immediately that, if $F$ is defined by the $A$-module $M$, then $F|U$ is defined by the $A_S$-module $M_S=A_S\otimes_A M$. The subsheaf $G$ of $F|U$ is then defined by a submodule $N'$ of $M_S$, and we see that $\overline G$ is precisely defined by the submodule $N$ of $M$ that is the inverse image of $N'$ under the canonical homomorphism from $M$ to $M_S$.
\begin{flushright}q.e.d.\end{flushright}

\sgaStableTarget{sga6.e00.cor.1}\textbf{Corollary 1 to proposition 2.10:} \emph{Let $x\in\underline K(X)$. For $x$ to have filtration $\geq i$, it is necessary and sufficient that there exist a closed subset $Y$ of $X$, of dimension $\leq n-i$, such that the restriction of $x$ to $X-Y$ is zero.}

%%%%%%%%%%  scan idx 71  |  folio 58  %%%%%%%%%%
\SGAPrintPage{58}%
\sgaStableTarget{sga6.e00.cor.2.occ2}\textbf{Corollary 2.} \emph{Assume $X$ connected of dimension $n$, without singularities in dimension $n-1$. If $D$ and $D'$ are two linearly equivalent divisors on $X$, we have $\gamma(D)\equiv\gamma(D')\quad\mathrm{mod.}\ \underline K(X)^2$.}

By corollary 1, we can assume $X$ nonsingular; the corollary then follows from the general formula:
\sgaStableTarget{sga6.e00.eq.2-31}\begin{equation}\tag{2.31}
 \gamma(D)=1-\gamma(L(-D))\qquad\mathrm{mod}\ \underline K(X)^2,
\end{equation}
where $L(-D)$ denotes the vector bundle of rank $1$ defined by the divisor $-D$. To verify (2.31), we write it in the form:
\sgaStableTarget{sga6.e00.eq.2-31bis}\begin{equation}\tag{2.31 bis}
 \gamma(L(-D))=1-\gamma(D)\qquad\mathrm{mod}\ \underline K(X)^2,
\end{equation}
and note that both members are homomorphisms from the group of divisors to the multiplicative group of invertible elements of the ring $\underline K(X)/\underline K(X)^2$ (note that $\underline K(X)^1/\underline K(X)^2$ has square zero, which is very easily verified from the explicit formula (2.15), using corollary 1). We are thus reduced to the case where $D$ is an irreducible hypersurface. More generally, if $D$ is a hypersurface whose components are disjoint, we have the equality:
\sgaStableTarget{sga6.e00.eq.2-32}\begin{equation}\tag{2.32}
 \gamma(D)=1-\gamma(L(-D))
\end{equation}
as follows from the exact sequence of sheaves
\[
0\longrightarrow\underline{O}_X(L(-D))\longrightarrow\underline{O}_X\longrightarrow\underline{O}_D\longrightarrow0.
\]

\sgaStableTarget{sga6.e00.thm.2-12}\textbf{Theorem 2.12.} \emph{Let $X$ be a nonsingular quasi-projective variety, and $Z$ and $Z'$ two homogeneous cycles of codimension $p$ on $X$ that are linearly equivalent. We then have:}
\sgaStableTarget{sga6.e00.eq.2-33}\begin{equation}\tag{2.33}
 \gamma_X(Z)\equiv\gamma_X(Z')\qquad\mathrm{mod}\ \underline K(X)^{p+1}.
\end{equation}
By hypothesis, we can find a cycle $Y$ of dimension $n-p+1$ on
%%%%%%%%%%  scan idx 72  |  folio 59  %%%%%%%%%%
\SGAPrintPage{59}%
$X\times k$ ($n$ is the
dimension of $X$), and two points $a,a'\in k$ such that we have:
\[
 Z=f_a^{-1}(Y),\qquad Z'=f_{a'}^{-1}(Y),
\]
where we put $f_u(x)=(x,u)$. By formula (2.23), we then have:
\[
 \gamma_X(Z)\equiv f_a^!(\gamma_{X\times k}(Y))\qquad\mathrm{mod}\ \underline K(X)^{p+1}
\]
and an analogous formula for $\gamma_X(Z')$. It then suffices to prove that $f_a^!=f_{a'}^!$, which follows immediately from proposition 2.9.
\begin{flushright}q.e.d.\end{flushright}

\sgaStableTarget{sga6.e00.cor.1.occ2}\textbf{Corollary 1:} \emph{The filtration of $\underline K(X)$ is compatible with its ring structure, i.e.\ we have $\underline K(X)^i\underline K(X)^j\subset\underline K(X)^{i+j}$ for all integers $i,j\geq0$. Moreover, we have a natural surjective homomorphism of graded rings:}
\sgaStableTarget{sga6.e00.eq.2-34}\begin{equation}\tag{2.34}
 \varphi:A(X)\longrightarrow G(\underline K(X))
\end{equation}
(where the second member denotes the graded ring associated with $\underline K(X)$ equipped with the filtration $(\underline K(X)^i)_{i\geq0}$) \emph{defined by the condition}:
\[
 \varphi(Z^{n-p})=\gamma(Z^{n-p})\qquad\mathrm{mod}\ \underline K(X)^{p+1}.
\]
This corollary follows immediately from theorem 2.12, propositions 2.6 and 2.7, and Chow's theorem 2.1, by induction on the integer $i+j$.

\sgaStableTarget{sga6.e00.cor.2.occ3}\textbf{Corollary 2:} \emph{Let $X$ and $Y$ be two nonsingular quasi-projective varieties and $f$ a morphism from $X$ to $Y$. Then the homomorphism $f^!$ from $\underline K(Y)$ to $\underline K(X)$ is compatible with the filtrations, and we have a commutative diagram:}
\[
\begin{tikzcd}[column sep=large,row sep=large]
G(\underline K(Y)) \arrow[r,"G(f^!)"] & G(\underline K(X))\\
A(Y) \arrow[r,"f^*"] \arrow[u,"\varphi_Y"'] & A(X) \arrow[u,"\varphi_X"'] .
\end{tikzcd}
\]
%%%%%%%%%%  scan idx 73  |  folio 60  %%%%%%%%%%
\SGAPrintPage{60}%
Using the graph of $f$, we decompose $f$ into the product of an \emph{immersion} of $X$ into $X\times Y$ and a \emph{projection} from $X\times Y$ onto $Y$. For the latter the corollary is trivial, and for an immersion we proceed as for corollary 1, using proposition 2.8 in place of proposition 2.7.

\sgaStableTarget{sga6.e00.prop.2-13}\textbf{Proposition 2.13 (``cellular decomposition''):} \emph{Let $P$ be an algebraic space (irreducible or not), equipped with an increasing family $P^0\subset P^1\subset\cdots\subset P^n=P$ of closed subsets such that, for every $i$ with $0\leq i\leq n$, the connected components of $P^i-P^{i-1}$ are algebraic spaces $V_{i,\alpha}$ isomorphic to $k^i$ (put $P^{-1}=\varnothing$). Then:}
\begin{enumerate}[label=\alph*)]
\item \emph{For every algebraic space $X$, the natural map from $\underline K(X)\otimes\underline K(P)$ to $\underline K(X\times P)$ deduced from the tensor product of sheaves on $X$ and $P$ is surjective.}
\item \emph{The $\gamma_P(V_{i,\alpha})$ form a system of generators of the group $\underline K(P)$.}
\end{enumerate}
The proof proceeds by induction on $n$, using the corollary of proposition 2.9 and proposition 2.10. (N.B. Unfortunately, I do not know whether the $\gamma_P(V_{i,\alpha})$ form a basis of $\underline K(P)$ and whether the homomorphism considered in a) is bijective. We shall not need to know this in what follows).

\sgaStableTarget{sga6.e00.thm.2-14}\textbf{Theorem 2.14.}
\emph{Let $x_i$ ($1\leq i\leq q$) and $E_j$ ($1\leq j\leq r$) be letters, and $n_j$ ($1\leq j\leq r$) integers $>0$; consider the symbols $\lambda^k(x_i)$ with $k\geq1$ and $\lambda^k(E_j)$ with $1\leq k\leq n_j$ as indeterminates, and let $R$ be a polynomial with integer coefficients in these indeterminates. Assume that, for every graded ring $A$, when in $R$ we substitute for the $x_i$ elements $\xi_i$ and for the $E_j$ elements $\eta_j$ of
